\documentclass[11pt]{article}

\usepackage[margin=0.72in]{geometry}
\usepackage{fontspec}
\usepackage{xeCJK}
\usepackage{booktabs}
\usepackage{longtable}
\usepackage{array}
\usepackage{xcolor}
\usepackage{hyperref}
\usepackage{fvextra}
\usepackage{enumitem}
\usepackage{amsmath}
\usepackage{titlesec}

\setmainfont{Arial}
\setCJKmainfont{Microsoft JhengHei}
\setCJKmonofont{Microsoft JhengHei}
\hypersetup{colorlinks=true, linkcolor=blue!60!black, urlcolor=blue!60!black}
\fvset{breaklines=true, breakanywhere=true, fontsize=\footnotesize}
\setlist[itemize]{leftmargin=1.35em}
\setlist[enumerate]{leftmargin=1.45em}
\titleformat{\section}{\Large\bfseries}{\thesection}{0.6em}{}
\titleformat{\subsection}{\large\bfseries}{\thesubsection}{0.6em}{}
\newcommand{\code}[1]{\texttt{\detokenize{#1}}}
\newcommand{\decision}[1]{\textbf{決策：}#1}
\newcommand{\impact}[1]{\textbf{影響：}#1}
\newcommand{\example}[1]{\textbf{例子：}#1}
\newcommand{\smallgap}{\vspace{0.35em}}

\title{\textbf{Optimization v2 決策目錄}\\
\large 實作細節、Prompt、Threshold、公式與 Grace/ICSI 範例}
\author{Virtual Mentor Memory System}
\date{2026-06-15}

\begin{document}
\maketitle

\section*{摘要}

這份文件取代早期的 \code{optimization_v2_decision_catalog_20260603}。早期版本已經描述第一版 v2，但仍偏英文，且沒有完整列出每個會影響系統結果的判定細節。這一版改成中文主體，目的很明確：讓教授或評審可以逐條審查系統如何做決定。

\smallgap
\noindent 本文件會回答：
\begin{itemize}
  \item 什麼內容會從 transcript 變成 L1 evidence object？
  \item 什麼 L1 會進入 L2？什麼 L1 會留在 unlinked / review-only？
  \item \code{related_topics} 到底有沒有用？它的權重多大？
  \item L2 label 怎麼選？公式是什麼？
  \item 什麼 L2 會升成 L3？L3 child 怎麼拆？何時會停止或進 review？
  \item Retrieval 是怎麼從 L1 找證據，再往上帶 L2/L3？
  \item 系統有哪些 prompt？原文是什麼？中文意思是什麼？
  \item Grace 和 ICSI 的具體例子如何對應到這些規則？
\end{itemize}

\smallgap
\noindent\textbf{最重要的邊界：} optimization v2 不直接改 raw L1 evidence。raw L1 來自 \code{share_mem/} 或 ICSI 的 \code{share_mem_effective/} sidecar。v2 只在 \code{optimization/runs/} 和 \code{optimization/reports/} 裡產生衍生 artifacts。這代表「dump」不是刪掉 L1，而是不讓它進入 active L2/L3 surface，或把它標成 unlinked、review-only、suppressed、warning。

\section{目前版本與參考輸出}

\begin{longtable}{p{0.28\linewidth}rr}
\toprule
\textbf{項目} & \textbf{Grace 最新 v2} & \textbf{ICSI BMR 29 場 effective}\\
\midrule
Run root & \multicolumn{2}{p{0.55\linewidth}}{\code{optimization/runs/grace_v2_latest_20260614}；\code{optimization/runs/icsi_bmr_full_completed29_v2_20260614}}\\
Profile & \code{mentor_mentee.yaml} & \code{isci_meeting.yaml}\\
Build mode & deterministic & deterministic\\
Build 時是否呼叫 LLM & 否 & 否\\
Source L1 數 & 448 & 6,233\\
Semantic key rows & 448 & 6,233\\
L2 topic 數 & 49 & 530\\
Linked L1 & 440 & 6,203\\
Unlinked L1 & 8 & 30\\
Raw materialized L3 parents & 3 & 65\\
Runtime L3 parents & 3 & 6\\
L3-assigned L1 & 71 & 2,683\\
L2 severe / warning & 0 / 2 & 0 / 9\\
L3 severe / warning & 0 / 2 & 0 / 9\\
\bottomrule
\end{longtable}

\noindent ICSI 的 raw L3 和 runtime L3 數量不同是刻意設計。Raw L3 用來審查所有可能的 parent/child split；runtime L3 是給 retrieval / demo 用的壓縮版 topic-family surface，避免把 65 個 L3 全部塞進 retrieval context。

\section{術語：Canonical、Optimization、Effective、Dump}

\begin{longtable}{p{0.22\linewidth}p{0.70\linewidth}}
\toprule
\textbf{術語} & \textbf{定義}\\
\midrule
Canonical & 目前既有主線 artifacts，例如 \code{share_mem/}、\code{long_term/l2/}、\code{long_term/l3/}。optimization v2 預設不覆蓋它們。\\
Optimization v2 & 新版 L2/L3 pipeline，位於 \code{optimization/long_term_v2/}。它讀 L1，產生新的 L2/L3 sidecar。\\
Raw L1 & 原始 evidence object。這是證據來源，不應被 L2/L3 pipeline 任意改寫。\\
Effective L1 & 經 sidecar gate 後的有效輸入。例如 ICSI 會使用 \code{filtered_share_mem} 或 \code{share_mem_effective}，避免 raw setup/noise 直接污染 L2/L3。\\
Unlinked & L1 保留在 evidence root，但沒有被分到 durable L2。\\
Review-only & 系統認為該項目可能有價值，但不夠安全自動放進 active L2/L3，需要人工審查。\\
Suppressed & topic review sidecar 決定某 L2 不出現在 effective topic surface。raw topic 還在 run root，可回查。\\
Dump & 本文口語上指「不進 active L2/L3」，不是刪除 raw L1。\\
\bottomrule
\end{longtable}

\section{整體資料流}

\begin{enumerate}
  \item Transcript 由上游 \code{share_mem/l1} pipeline 轉成 L1 evidence objects。
  \item 若是 ICSI 這種 noisy corpus，先用 sidecar gate 形成 effective L1 root。
  \item optimization v2 讀取 L1 root，複製 input snapshot 並計算 hash。
  \item 每個 L1 產生 semantic key row。
  \item L2 induction 根據 semantic keys、content/evidence、importance、optional related topics 形成 durable topics。
  \item 過大的 L2 才考慮升成 L3 topic family。
  \item Runtime export 使用 effective topic surface：topic review sidecar、L3 child review sidecar、corpus-theme L3 sidecar 會影響最後 retrieval 看到的 L2/L3。
  \item Retrieval 先找 L1 evidence seeds，再帶出 L2/L3 context。
\end{enumerate}

\section{L1：什麼會被抽成 evidence object}

Optimization v2 不重跑 L1，但 L1 是整個系統的證據基礎，所以必須說清楚 upstream L1 的決策。

\subsection{上游 L1 pipeline 的固定參數}

\begin{longtable}{p{0.32\linewidth}p{0.22\linewidth}p{0.38\linewidth}}
\toprule
\textbf{決策點} & \textbf{Threshold / 值} & \textbf{意義}\\
\midrule
Grace chunk size & 24 transcript lines & Grace 中文 turn 較長，一個 line 常包含很多上下文，所以視窗較小。\\
ICSI chunk size & 40 transcript lines & ICSI 英文 speaker turns 較短且碎，因此視窗較大。\\
Issue episode gap & Grace 10；ICSI 15 lines & 同一問題相隔不遠時可以被視為同一 episode。\\
Segmentation target & 通常 3--6 segments & segment 是 topic-coherent discussion，不是逐句切。\\
Max segments per window & 10 & 只有明顯大主題切換時才允許這麼多。\\
Preferred segment length & 10--24 primary-window lines & 防止過度切碎，也防止一段過大。\\
Idea unit target & 通常 3--8 units / segment & 每個 unit 是一個可檢查 idea。\\
Max idea units per segment & 12 & 只有一段內有很多 distinct durable claims 時使用。\\
Extraction batch target & 10 idea units & typed L1 agents 讀 bounded packet。\\
Max extraction batch & 14 idea units & 超過會切 batch。\\
Batch overlap & 2 idea units & 避免 batch 邊界漏掉上下文。\\
Typed L1 candidates & 通常 0--4，最多 6 & 避免把每個 idea unit 都重述成 L1。\\
Fallback candidates & 最多 3 & typed agents 無候選時才啟用。\\
Previous context & 前 2 batches，最多 8 items & 只能消歧，不是 evidence。\\
Previous item chars & 160 chars & 避免 previous context 壓過 bounded evidence。\\
Cross-window topic merge & Jaccard $\geq 0.30$ & topic label 相似時觸發跨 window refinement。\\
Cross-window boundary merge & Jaccard $\geq 0.20$ & 邊界 transcript 相似時觸發 refinement。\\
Cross-window idea-unit merge & Jaccard $\geq 0.20$ & idea unit 相似時觸發 refinement。\\
Boundary refinement context & 16 lines lookback/lookahead & 重切跨 window 邊界時讀更多上下文。\\
\bottomrule
\end{longtable}

\subsection{什麼內容會成為 L1}

\decision{只有 durable long-term memory 會被提成 L1。}

\begin{itemize}
  \item 會保留：專案級 decision、method change、concrete follow-up、stable finding、unresolved research question、decision-supporting argument。
  \item 會保留：會影響未來系統行為的 requirement / scope clarification。
  \item 會保留：evaluation design、baseline、metric、full-context/RAG/ablation comparison 條件。
  \item 會保留：模型能力判斷，若它影響 pipeline design 或 evaluation。
  \item 會丟到非 L1 或低重要度：filler、acknowledgement、local wording clarification、社交話語、一句話 example、單純 setup/logistics。
  \item 重要度規則：\(\geq 0.90\) 只給 project-level 或 future-steering items；\(0.50\)--\(0.70\) 用於有用但偏 local meeting context。
\end{itemize}

\subsection{L1 例子}

\example{Grace L1-0422-006}
\begin{quote}
為了取代逐「輪」更新記憶體的簡單作法，團隊探討了兩種新的方法來預處理逐字稿以生成 idea units：迭代法與預分割法。
\end{quote}
\textbf{為什麼成為 L1：}這是 memory update pipeline 的方法改變，會影響後續設計，不是單次聊天內容。

\smallgap
\example{Grace L1-0307-034}
\begin{quote}
團隊成員需要協調統一資料檔案夾結構，以避免檔案重複。
\end{quote}
\textbf{為什麼可能不進 L2：}它是有效 L1，但在目前 corpus 中不是 recurring durable topic。v2 保留 evidence，但不硬塞進 L2。

\smallgap
\example{ICSI L1-Bmr002-094}
\begin{quote}
The project needs to select an annotation tool. Tools from Mississippi State and XWaves have been suggested as potential options.
\end{quote}
\textbf{為什麼成為 L1：}這是 annotation workflow 的工具選擇問題，對後續 corpus annotation 有持續影響。

\section{ICSI Effective L1 Gate：為什麼不直接讀 raw share\_mem}

\decision{ICSI L2/L3 優先讀 \code{filtered_share_mem} 或 \code{share_mem_effective}，不直接把 raw L1 全部視為 active evidence。}

\begin{itemize}
  \item raw L1 仍保留，方便審查。
  \item effective L1 移除或降權 source-data-only、same-evidence duplicate、microphone inventory、kickoff metadata、speaker-code/filler 等項目。
  \item 這不是為 ICSI 硬編答案，而是 mentor-mentee / research meeting 場景常見的 evidence hygiene：source data 和 durable memory 要分開。
\end{itemize}

\section{Semantic Key Layer：L1 如何變成可分群訊號}

\subsection{輸入與輸出}

每個 L1 object 會輸入：
\[
content,\ evidence,\ type,\ importance,\ related\_topics,\ related\_obj\_ids
\]

輸出 semantic key row：
\begin{itemize}
  \item \code{candidate_terms}
  \item \code{semantic_facets}
  \item \code{source_signals}
  \item \code{confidence}
\end{itemize}

\subsection{Term candidate 的產生}

\decision{v2 從 content/evidence 產生 1--3 gram terms，再用 profile stopwords/generic labels/type-like labels/reject labels 過濾。}

若 profile 的 \code{topic_key_language} 是 English，且有英文 candidate，則排除 CJK candidate，避免 ICSI/Grace mixed text 把語言殘片當 topic key。

\subsection{Term ranking 公式}

對每個 candidate term：
\[
\begin{aligned}
score &= 2.5 \cdot \min(count,4) + phrase\_bonus
      + 0.25 \cdot \min(specific\_token\_count,3)\\
score &\mathrel{-}= 0.9 \quad \text{if single token}\\
score &\mathrel{-}= 2.5 \quad \text{if single token and generic}\\
score &\mathrel{-}= 0.4 \quad \text{if multi-token but all generic}\\
score &\mathrel{-}= 0.15 \cdot generic\_token\_count
\end{aligned}
\]

\[
phrase\_bonus =
\begin{cases}
1.6 & \text{三個以上 tokens}\\
1.15 & \text{兩個 tokens}\\
0 & \text{其他}
\end{cases}
\]

排序時使用 \(-score\)、phrase 長度、字母序作 deterministic tie-break。

\subsection{Semantic key confidence}

\[
confidence = \min(1.0,\ 0.25 + 0.08 \cdot \#facet\_terms)
\]

\impact{這個 confidence 不代表最終 topic 正確，只代表該 L1 產生的 semantic signals 是否足夠。}

\section{L2：L1 什麼時候會被 assign}

\subsection{Profile thresholds}

\begin{longtable}{p{0.35\linewidth}p{0.25\linewidth}p{0.32\linewidth}}
\toprule
\textbf{參數} & \textbf{Grace / mentor\_mentee} & \textbf{ICSI}\\
\midrule
\code{min_topic_evidence} & 4 & 繼承 4\\
\code{minimum_assignment_confidence} & 0.25 & 繼承 0.25\\
\code{high_importance_threshold} & 0.70 & 繼承 0.70\\
\code{max_l2_topic_share_before_specificity_penalty} & 0.18 & 繼承 0.18\\
\code{allow_single_meeting_topic_if_high_importance} & true & false\\
\code{review_only_min_assignment_confidence} & 0.25 default & 0.52\\
\code{semantic_rescue_min_confidence} & 0.60 default & 0.68\\
\code{max_review_only_l1_share} & 無硬限制 & 0.20\\
\bottomrule
\end{longtable}

\subsection{Label selection 公式}

對 L1 的每個 candidate label，令 \(df\) 為 document frequency、\(N\) 為 L1 總數、\(m\) 為 \code{min_topic_evidence}、\(max\_share=0.18\)：

\[
\begin{aligned}
capped\_df &= \min(df, 3m)\\
share &= \frac{df}{\max(N,1)}\\
score &= 1.4 \cdot capped\_df + phrase\_bonus
       + 0.2 \cdot \min(specific\_token\_count,3)\\
score &\mathrel{+}= 2.0 \quad \text{if } df \geq m\\
score &\mathrel{-}= 30.0 \cdot (share - max\_share) \quad \text{if } share > max\_share\\
score &\mathrel{-}= 3.0 \quad \text{if single token}\\
score &\mathrel{+}= 1.0 \quad \text{if from content or facet}\\
score &\mathrel{+}= \max(0,\ 2.0 - 0.25\cdot rank) \quad \text{if content/facet phrase has two or more tokens}\\
score &\mathrel{+}= 0.25 \quad \text{if only from related topics}
\end{aligned}
\]

\[
phrase\_bonus =
\begin{cases}
1.15 & \text{兩個 tokens}\\
0.65 & \text{三個以上 tokens}\\
0 & \text{其他}
\end{cases}
\]

\decision{\code{related_topics} 有用，但不是主因。} content/facet signal 至少加 \(1.0\)，而 related topic only 只加 \(0.25\)。這避免新資料集被舊 topic vocabulary 強迫拉歪。

\subsection{Assignment confidence 公式}

\[
confidence = \min\left(0.95,\ 0.2 + 0.35\cdot\frac{df}{m} + 0.25\cdot importance\right)
\]

若是 semantic rescue，使用：
\[
rescue\_score = 0.7 \cdot containment + 0.3 \cdot Jaccard(row\_signature,\ topic\_signature)
\]
\[
containment = \frac{|shared\_terms|}{\min(|row\_signature|,\ |topic\_signature|)}
\]
\[
rescue\_confidence = \min(0.82,\ 0.32 + 0.6\cdot rescue\_score + 0.12\cdot importance)
\]

Rescue 必須至少有兩個 shared terms，且 \(rescue\_score \geq 0.18\)。

\subsection{Assign L2 的條件}

L1 會被 assign 到 L2，通常需要同時滿足：
\begin{enumerate}
  \item 有通過 filter 的 candidate label。
  \item label 的 \(df \geq min\_topic\_evidence\)，或該 L1 importance 足夠高且 profile 允許 single/high-importance topic。
  \item assignment confidence \(\geq minimum\_assignment\_confidence\)。
  \item label 不是 type-like、generic-only、reject label、weak phrase。
  \item 如果是 singleton topic，profile 允許，否則轉 unlinked/review。
\end{enumerate}

\subsection{不進 L2 / dump 的條件}

\begin{longtable}{p{0.30\linewidth}p{0.26\linewidth}p{0.36\linewidth}}
\toprule
\textbf{原因碼} & \textbf{觸發條件} & \textbf{結果}\\
\midrule
\code{no_repeated_or_high_confidence_semantic_key} & 沒有可用 label & L1 保留，但 unlinked。\\
\code{candidate_topic_not_durable_enough} & \(df < m\) 且 importance 未達 high threshold & 不建立不耐久 L2。\\
\code{low_assignment_confidence} & confidence 低於 minimum & unlinked 或 review。\\
\code{review_only_low_assignment_confidence} & 有 recurring label 但低於 review-only 門檻 & 進 manual review，不進 active L2。\\
\code{high_importance_ambiguous_topic_review} & 高 importance 但 topic ambiguous & 保留 evidence，要求人工判斷。\\
\code{singleton_l2_not_durable_enough} & 該 L2 只有一個 L1，且 profile 不允許 & 不建立 singleton L2。\\
\code{semantic_rescue_low_confidence_review} & rescue 有 overlap 但 confidence 不足 & review，不直接 assign。\\
\bottomrule
\end{longtable}

\section{L2 Validation 與 Label Polish}

\subsection{Validation severe 條件}

\begin{itemize}
  \item 缺少必要輸出：\code{manifest.json}、input snapshot、semantic key index、L2/L3 view/index。
  \item L2 topic 沒有 \code{representative_l1_ids}。
  \item L2 topic 沒有 \code{creation_rationale}。
  \item L2 assignment 沒有 \code{rationale}。
\end{itemize}

\subsection{Validation warning 條件}

\begin{itemize}
  \item L1 text 似乎有 mojibake / private-use glyph。
  \item semantic key count 和 source L1 count 不一致。
  \item L2 label 是 type-like、generic、rejected、weak phrase。
  \item L2 只有一個 linked L1。
  \item assignment 用到 related topics，但 rationale 太弱。
\end{itemize}

\subsection{Label polish sidecar}

Label polish 不直接改 raw L2/L3，而是產生 review queue。觸發條件：

\begin{longtable}{p{0.28\linewidth}p{0.30\linewidth}p{0.34\linewidth}}
\toprule
\textbf{類別} & \textbf{例子 / 規則} & \textbf{結果}\\
\midrule
Type-like label & decision, proposal, argument, finding, action item, open issue & label review。\\
Generic-only label & data, system, memory, discussion, topic, thing, process & label review。\\
Single artifact label & uh, um, hmm, okay, yeah, go, ahead & suppress/review。\\
Weak terminal word & label 結尾是 of, in, to, with, per 等 & label review。\\
Weak start word & basically, maybe, probably, actually, just 等 & label review。\\
Ambiguous short label & 長度 \(\leq 4\) 且不是全大寫 & label review。\\
Missing representative L1 & label 沒有代表 evidence & severe/review。\\
\bottomrule
\end{longtable}

\section{L3：L2 什麼時候升成 topic family}

\subsection{Promotion threshold}

L3 是 topic-family navigation，不是主要 evidence。只有 oversized L2 才會升。

令 \(sizes\) 是所有 L2 的 linked-L1 counts：
\[
p85 = percentile(sizes, 0.85)
\]

若總 L1 \(\leq 50\)：
\[
threshold = \max(min\_absolute\_threshold,\ p85)
\]

若總 L1 \(> 50\)：
\[
threshold = \max(min\_absolute\_threshold,\ target\_child\_max,\ p85)
\]

目前校準值：
\begin{itemize}
  \item Grace: \code{l3_absolute_l1_threshold = 49}，target child size 6--20。
  \item ICSI: \code{l3_absolute_l1_threshold = 131}，target child size 6--20。
\end{itemize}

\subsection{Target child count}

\[
adaptive\_child\_count = \left\lceil \frac{event\_count}{target\_child\_max} \right\rceil
\]
\[
child\_count = clamp(adaptive\_child\_count,\ min\_child\_count,\ max\_child\_count)
\]

Default: \code{min_child_l2_count = 2}，\code{max_child_l2_count = 12}。

\subsection{Child label 怎麼產生}

\decision{v2 不使用舊版手寫 child taxonomy。}

Child labels 來自：
\begin{itemize}
  \item parent L2 的 top semantic terms；
  \item parent timeline summaries；
  \item 2--3 gram candidates；
  \item profile label filters / normalization aliases。
\end{itemize}

排序：
\[
rank\_score = count + 4 \cdot token\_count
\]
再用 phrase 長度與字母序 tie-break。

\subsection{Event assign child 的規則}

每個 parent L2 timeline event 會被分到最相符 child：
\[
score(event, child) = |tokens(event.summary) \cap tokens(child.assignment\_criteria)|
\]

若所有 child 分數都 \(\leq 0\)，用 deterministic round-robin fallback：
\[
child = candidates[index \bmod child\_count]
\]

\impact{這個 fallback 不是表示 child split 品質好；後續 separability validation 會擋掉低品質 split 或放進 review。}

\subsection{Split separability score}

Split 初始分數 \(1.0\)，再扣分：

\begin{longtable}{p{0.45\linewidth}p{0.16\linewidth}p{0.31\linewidth}}
\toprule
\textbf{條件} & \textbf{扣分} & \textbf{原因碼}\\
\midrule
Non-empty children 少於 2 個 & 0.60 & \code{insufficient_non_empty_children}\\
最大 child 佔比 \(\geq 0.80\) 且超過 max child size & 0.25 & \code{dominant_oversized_child}\\
超過一半 children 為空，且最大 child 超大 & 0.15 & \code{too_many_empty_children}\\
Sibling label overlap \(\geq 0.72\) & 0.25 & \code{high_sibling_label_overlap}\\
Weak child labels & 最多 0.25 & \code{weak_child_labels}\\
Manual-review child labels & 最多 0.25 & \code{manual_review_child_labels}\\
\bottomrule
\end{longtable}

Minimum split separability:
\begin{itemize}
  \item mentor-mentee default: 0.35。
  \item ICSI: 0.58。
\end{itemize}

\subsection{L3 不 materialize 的條件}

\begin{itemize}
  \item L2 event count 小於 adaptive threshold。
  \item child candidates 不足，或 split quality 太差。
  \item 大部分 evidence 仍落在同一 child，其他 child 只是 tiny edge cases。
  \item child labels 是 generic/type-like/weak/rejected。
  \item split separability 低於 profile threshold。
\end{itemize}

\section{Effective Topic Surface：最後 retrieval 看到什麼}

Runtime 不是永遠直接讀 raw \code{l2_view.json} / \code{l3_view.json}。載入順序：

\begin{enumerate}
  \item L2：如果 \code{topic_review/effective_topic_index.json} 有 suppressed L2 ids，active L2 view 會過濾它們。
  \item L3：若存在 \code{l3/corpus_theme_view.json} 和 \code{l3/corpus_theme_index.json}，優先使用 corpus-theme L3。
  \item 若無 corpus-theme L3，但有 \code{effective_l3_view.json}，使用 L3 child review 後的 effective L3。
  \item 否則使用 raw L3。
\end{enumerate}

優先序：
\[
corpus\_theme\_l3 > l3\_child\_review > raw\_l3
\]

\section{Retrieval：如何從 L1 往上帶 L2/L3}

\subsection{L1 ranking 公式}

Retrieval 先找 L1 evidence seeds。對 query \(q\)、L1 text \(t\)、topic labels \(r\)：
\[
base = \max(Jaccard(q,t),\ 0.95\cdot Jaccard(q,r))
\]

Phrase bonus:
\[
phrase\_bonus += 0.14 \quad \text{if topic phrase appears in query phrase}
\]
\[
phrase\_bonus += 0.08 \quad \text{if topic tokens are subset of query tokens}
\]

Final score:
\[
score = base + \min(0.35,\ phrase\_bonus) + 0.015\cdot importance
\]

分數 \(\leq 0\) 的 L1 不會被選入候選。

\subsection{Context construction}

目前 ICSI comparison 使用：
\begin{itemize}
  \item RAG: top-k L1 = 20。
  \item Optimization v2 layered: top-k L1 = 60。
  \item Full context: 幾乎所有 L1，平均 4,715 L1。
\end{itemize}

Layered context 不是只列 L1，也會附：
\begin{itemize}
  \item selected L1 evidence seeds；
  \item selected L2 labels / timelines；
  \item selected L3 labels；
  \item evidence grounding metadata。
\end{itemize}

\section{ICSI 評估結果}

目前 ICSI 29 場使用 20 題 held-out evidence questions，這是 no-LLM retrieval/evidence comparison，不是最終 answer benchmark。

\begin{longtable}{p{0.30\linewidth}rrr}
\toprule
\textbf{系統} & \textbf{Expected L1 recall} & \textbf{Avg context tokens} & \textbf{Avg selected L1}\\
\midrule
Full context L1 & 1.0000 & 575,199 & 4,715\\
Traditional lexical RAG & 0.5583 & 2,369.95 & 20\\
Optimization v2 layered & 0.8958 & 7,358.35 & 60\\
\bottomrule
\end{longtable}

\noindent\textbf{解讀：}
Full context recall 最高，因為它幾乎全塞；但 token 成本極高。RAG 很省，但漏掉很多 expected L1。Optimization v2 token 比 RAG 多，但 recall 明顯高，而且提供 L2/L3 context。

\section{具體範例：Grace}

\subsection{問題：為什麼要把 transcript 拆成 segment / idea units？}

會被取出的關鍵 evidence：
\begin{itemize}
  \item \textbf{L1-0422-006}：為了取代逐 turn update，討論 idea units 的迭代法與預分割法。
  \item \textbf{L1-0429-122}：放棄 full transcript 一次更新，因為 context 太大且 JSON output 難控制。
  \item \textbf{L1-0506-001}：確立 context planner \(\rightarrow\) segment agent \(\rightarrow\) idea units。
  \item \textbf{L1-0506-002}：定義 segment 和 idea unit 差異。
\end{itemize}

\textbf{L2/L3 的角色：}RAG 可能找到其中一段，但 L2/L3 可以把跨 0422、0429、0506 的演進串起來，說明從 full transcript / turn-based update 到 segment/idea-unit pipeline 的原因。

\section{具體範例：ICSI}

\subsection{Annotation tool}

\begin{quote}
\textbf{L1-Bmr002-094}：The project needs to select an annotation tool. Tools from Mississippi State and XWaves have been suggested.
\end{quote}

L2: \code{annotation tool}。Runtime L3: annotation and transcription workflow。

\textbf{判斷原因：}這不是 Grace memory-system topic，而是 ICSI corpus annotation workflow 的 recurring evidence-backed topic。

\subsection{Disk space}

\begin{quote}
\textbf{L1-Bmr024-097}：A decision was made to save all recorded data, including distant microphones, because disk space is inexpensive.
\end{quote}

L2: \code{disk space}。L3: corpus design and data management。

\textbf{判斷原因：}這是 corpus storage policy，不是單純硬體聊天。

\subsection{Resource allocation}

\begin{quote}
\textbf{L1-Bmr001-052}：Speaker location detection is feasible with available data, but who will write the software is unresolved.
\end{quote}

L2: \code{resource allocation}。L3: project planning and management。

\textbf{判斷原因：}這是 unresolved implementation ownership，會影響後續 project planning。

\section{Prompt 完整清單}

\subsection{說明}

Optimization v2 deterministic build 不呼叫 LLM。以下 prompt 分兩類：
\begin{itemize}
  \item 上游 L1 extraction prompt：用於 transcript \(\rightarrow\) L1。
  \item Optimization v2 optional / diagnostic prompt：用於 label refinement、split review、L3 child review、answer generation、answer judging。
\end{itemize}

因為程式實際送出的 prompt 是英文，所以本節保留英文原文摘要，並附中文翻譯。

\subsection{P-L1-1 Segmentation Agent}

\textbf{原文摘要}
\begin{Verbatim}
You are segmentation_agent for a long-term memory extraction pipeline.
Return JSON only. Split the transcript window into topic-coherent segments.
Use original line numbers.

Create durable discussion segments, not sentence-level or checklist-like slices.
Normally return 3-6 segments for this primary window.
You may return up to 10 segments only when there are clear major topic shifts.
Prefer 10-24 primary-window lines per segment.
Do not split every small subpoint into a segment; downstream idea units will
handle the smaller claims inside each durable segment.
\end{Verbatim}

\textbf{中文翻譯}：你是 segmentation agent。請只回 JSON，把 transcript window 切成 topic-coherent segments，保留原始行號。Segment 要是可長期使用的討論段，不是句子級或 checklist 級切片。一般 3--6 段，只有明顯大主題切換時最多 10 段，每段最好 10--24 行。

\subsection{P-L1-2 Idea Unit Agent}

\textbf{原文摘要}
\begin{Verbatim}
You are idea_unit_agent. Convert this segment into compact idea units.
Each unit should express one checkable idea that downstream L1 agents can share.
Do not classify memory types here. Return JSON only.

Normally return 3-8 idea units. You may return up to 12 only when the segment
contains many distinct durable claims.
Do not split every sentence into a separate unit.
Preserve questions, objections, counterexamples, method exploration, evaluation
criteria, comparison rationale, and unresolved design discussions when they
affect future project decisions or memory behavior.
\end{Verbatim}

\textbf{中文翻譯}：你是 idea unit agent。請把 segment 轉成 compact idea units。每個 unit 要是一個可檢查、可讓下游 L1 agents 共用的 idea。此處不分類 memory type。一般 3--8 個，只有 segment 中有很多不同 durable claims 時最多 12 個。不要逐句切；要保留會影響未來決策或記憶行為的問題、反對、反例、方法探索、評估標準、比較理由與未解設計討論。

\subsection{P-L1-3 Idea Unit Repair Agent}

\textbf{原文摘要}
\begin{Verbatim}
You are idea_unit_repair_agent. Repair idea-unit coverage for one segment.
Return a full revised set of semantic idea units for the segment, not a patch.
Return JSON only.

The validator found issues in the first idea-unit pass. Your job is to avoid
deterministic fallback by creating durable semantic units where there is useful
project content, or marking true filler as non_memory_context_ranges.
\end{Verbatim}

\textbf{中文翻譯}：你是 idea unit repair agent。請修復某 segment 的 idea-unit coverage，回傳完整新版，不是 patch。validator 發現第一輪有問題，你的任務是在有用內容中建立 durable semantic units，或把真正 filler 標成 non-memory ranges。

\subsection{P-L1-4 Typed L1 Agent}

\textbf{原文摘要}
\begin{Verbatim}
You are l1_{obj_type}_agent in a multi-agent long-term memory pipeline.
Read only the bounded idea units below and propose only {obj_type} candidates.
Do not infer from outside this extraction scope.
Prior context may help disambiguate references, but it is never evidence.
Every candidate must be supported by the bounded idea units below.

Normally return 0-4 candidates. You may return up to 6 only when there are
clearly distinct durable {obj_type} memories.

Only output durable long-term memory:
- keep project-level decisions, method changes, concrete follow-ups, stable
  findings, unresolved research questions, or decision-supporting arguments
- do not restate each idea unit as a candidate
- do not output local clarifications, filler, examples, or one-line observations
- use importance >= 0.90 only for project-level or future-steering items
- use 0.50-0.70 for useful but local meeting-level context
Return JSON only.
\end{Verbatim}

\textbf{中文翻譯}：你是某個 L1 類型的 agent。只能讀 bounded idea units，只提出該類型候選，不得從 scope 外推論。prior context 只能消歧，不能作 evidence。一般回傳 0--4 個，只有明確不同的 durable memories 時最多 6 個。只輸出長期有用的記憶；不要把每個 idea unit 重述成 candidate。

\subsection{P-L1-5 General L1 Fallback Agent}

\textbf{原文摘要}
\begin{Verbatim}
You are general_l1_fallback_agent in a multi-agent long-term memory pipeline.
This fallback runs only because the typed L1 agents produced no candidates.

Read only the bounded idea units below. Propose a small number of durable L1
candidates only if this extraction packet clearly contains one of the allowed
types. Return an empty candidates list if the packet is purely filler,
logistics, acknowledgements, or unclear.
Return at most 3 candidates. Return JSON only.
\end{Verbatim}

\textbf{中文翻譯}：你是 general fallback agent。只有 typed L1 agents 沒有產生候選時才啟用。只讀 bounded idea units；只有 packet 明確包含允許的 durable L1 類型時才提出少量候選。若只是 filler、行政流程、附和或不清楚，回傳空清單。最多 3 個。

\subsection{P-V2-1 L2 Label Refinement}

\textbf{原文摘要}
\begin{Verbatim}
Task: Refine evidence-derived L2 topic labels for a mentor-mentee research
meeting memory system.

Constraints:
- Use only the provided L1 evidence summaries and semantic terms.
- Do not invent project-specific labels that are not supported by evidence.
- Prefer concise English noun phrases, usually 2 to 5 words.
- Prefer stable noun labels over verb phrases.
- Avoid labels that are merely L1 object types, generic buckets, or
  implementation artifacts.
- Keep a topic broad enough to cover its linked evidence but specific enough
  to be useful for retrieval.
\end{Verbatim}

\textbf{中文翻譯}：請改善 evidence-derived L2 labels。只能使用提供的 L1 summaries 和 semantic terms；不能發明 evidence 不支持的 label；偏好 2--5 詞英文名詞片語；避免 type-like、generic bucket 或 implementation artifact；label 要能覆蓋 linked evidence，但仍具體到有助 retrieval。

\subsection{P-V2-2 L2 Induction Review Proposal}

\textbf{原文摘要}
\begin{Verbatim}
Task: Review evidence-induced L2 topics and propose merge, split, or
assignment-review candidates. Do not rewrite the source data.

Constraints:
- Use only topic labels, semantic terms, and representative L1 evidence.
- Every proposal must cite representative_l1_ids.
- Prefer no proposal over speculative proposals.
- Do not propose Grace-specific taxonomy unless the evidence supports it.
- Return JSON only.
\end{Verbatim}

\textbf{中文翻譯}：請審查 evidence-induced L2 topics，提出 merge、split 或 assignment-review 候選，但不能改 source data。每個 proposal 必須引用 representative L1 ids。寧可沒有 proposal，也不要猜測。除非 evidence 支持，不得提出 Grace-specific taxonomy。

\subsection{P-V2-3 Focused Split Review}

\textbf{原文摘要}
\begin{Verbatim}
Task: For each oversized L2 that deterministic splitting could not safely
materialize, propose evidence-backed child L2 labels or say review_only.

Constraints:
- Use only source_l2_id, semantic terms, and timeline_sample evidence.
- Every split proposal must cite representative_l1_ids.
- Every child must include assignment_criteria and representative_l1_ids that
  distinguish it from sibling children.
- For large sources, propose enough child candidates to avoid one dominant
  residual bucket; if evidence does not support separable coverage, use
  review_only.
- Do not invent unsupported project-specific taxonomy.
\end{Verbatim}

\textbf{中文翻譯}：對 deterministic split 無法安全 materialize 的 oversized L2，請提出 evidence-backed child labels 或標成 review-only。每個 child 必須有 assignment criteria 和 representative L1 ids，能和 sibling 區分。若 evidence 不支持可分離覆蓋，就不要硬拆。

\subsection{P-V2-4 L3 Child Label Review}

\textbf{原文摘要}
\begin{Verbatim}
Task: Review L3 child topics for an evidence-grounded long-term memory system.

Instructions:
- Use only the provided representative L1 content.
- Do not invent a child topic that is not supported by the L1 evidence.
- Judge whether each child is coherent and whether the label names the evidence.
- Prefer concise English noun phrases for relabel decisions.
- If evidence is mixed or the child is not coherent, choose reject_split or
  needs_human_review.
\end{Verbatim}

\textbf{中文翻譯}：請審查 L3 child topics。只能使用提供的 representative L1 content。不能創造 evidence 不支持的 child topic。判斷 child 是否 coherent，以及 label 是否真的命名 evidence。若 evidence 混雜或 child 不 coherent，選 reject split 或 needs human review。

\subsection{P-V2-5 Answer Generation}

\textbf{原文}
\begin{Verbatim}
You are a meeting-memory QA assistant. Answer only from the provided context.
If the context is insufficient, say what is missing. Cite meeting ids, L1 obj_ids,
or transcript line ranges whenever possible. Keep the answer concise and factual.

Context:
{retrieved_context}

Question:
{query}

Answer:
\end{Verbatim}

\textbf{中文翻譯}：你是 meeting-memory QA assistant。只能從 provided context 回答；如果 context 不足，說明缺什麼。可能時引用 meeting id、L1 obj id 或 transcript line range。回答要簡潔、事實性。

\subsection{P-V2-6 Answer Judge}

\textbf{原文摘要}
\begin{Verbatim}
Use the expected answer as the primary gold standard. Use the retrieved context
as supplementary grounding and traceability evidence. The retrieved context may
be truncated for very large strategies such as full transcript, so do not mark
factual_correctness as zero merely because a cited detail is absent from the
excerpt.

Dimensions:
- factual_correctness
- evidence_grounding
- topic_evolution
- completeness
- conciseness
- hallucination_risk
- source_traceability
\end{Verbatim}

\textbf{中文翻譯}：expected answer 是主要 gold standard；retrieved context 是 grounding/traceability 的補充。full transcript 類策略 context 可能被截斷，所以不能只因 excerpt 沒看到細節就把 factual correctness 歸零。評分維度包含事實正確、證據支撐、topic evolution、完整性、簡潔性、幻覺風險、來源可追溯性。

\subsection{P-V2-7 ICSI Answer Judge}

\textbf{原文摘要}
\begin{Verbatim}
Use the expected gold evidence as the main target. Use the retrieved context
excerpt only to assess grounding and traceability. The full-context strategy may
have a very large retrieved context, so the excerpt can be incomplete; do not
set factual_correctness to zero merely because the answer cites a detail outside
the excerpt.

topic_evolution: high when the answer captures evolution over time when the
query asks for it; otherwise score neutral to high if not relevant.
hallucination_risk: 0.0 means no unsupported fabrication risk; 1.0 means severe
unsupported fabrication.
\end{Verbatim}

\textbf{中文翻譯}：ICSI 評分以 expected gold evidence 為主要目標。retrieved context excerpt 只用來輔助判斷 grounding 和 traceability。full context 可能太大而 excerpt 不完整，因此不能因 excerpt 缺某細節就把 factual correctness 歸零。topic evolution 只在問題要求演進時才應高要求；hallucination risk 是反向分數，0 代表沒有 unsupported fabrication。

\section{Reproduction Commands}

\begin{Verbatim}
uv run python optimization/long_term_v2/build_view.py `
  --share-mem-root share_mem `
  --profile optimization/long_term_v2/profiles/mentor_mentee.yaml `
  --out optimization/runs/grace_v2_latest_20260614 `
  --mode deterministic `
  --clean

uv run python optimization/long_term_v2/build_view.py `
  --share-mem-root memory_outputs/icsi/runs/bmr_full_completed29_eval_20260614/share_mem_effective `
  --profile optimization/long_term_v2/profiles/isci_meeting.yaml `
  --out optimization/runs/icsi_bmr_full_completed29_v2_20260614 `
  --mode deterministic `
  --clean

uv run python optimization/long_term_v2/validate_view.py `
  --run-root optimization/runs/icsi_bmr_full_completed29_v2_20260614
\end{Verbatim}

\section{結論}

Optimization v2 目前可以清楚解釋每個主要決策：L1 是 evidence；semantic keys 是中介訊號；L2 是 durable topic；L3 是 oversized topic 的 navigation layer；retrieval 先找 L1，再帶 L2/L3。最重要的是，v2 不再依賴 Grace-specific hand-coded taxonomy，也不讓 \code{related_topics} 成為唯一 assignment source。剩餘工作主要是 label polish、split review、以及 runtime opt-in 的穩定驗證，而不是重新大改架構。

\end{document}
